\chapter{Datos, conjunto de datos y supervisión débil}
\label{cap:7}

Este capítulo detalla la preparación del conjunto de datos y la metodología de supervisión débil utilizada para suplir la falta de etiquetas de prioridad históricas. En primer lugar, se expone el origen de los datos procedentes del Registro 112 de Castilla y León y el proceso de depuración aplicado a las variables. A continuación, se describe la lógica de supervisión débil que permite generar etiquetas académicas a partir de múltiples heurísticas operativas y se evalúa el acuerdo inter-etiquetador. Finalmente, se detalla la estrategia de partición de datos y los controles implementados para evitar el filtrado de información y asegurar la solidez de la evaluación.

El capítulo anterior ha definido la arquitectura del sistema. Este capítulo desarrolla la base empírica que permite alimentar, entrenar y evaluar esa arquitectura. En particular, explica el uso del Registro de Emergencias del 112 de Castilla y León, las variables V01--V15, la construcción de la etiqueta académica P1--P4, la supervisión débil, la prevención de fugas de información y las particiones de entrenamiento, validación y prueba.

El proyecto trabaja con registros reales, pero no dispone de una prioridad oficial. Por ello, el conjunto no puede tratarse como una fuente supervisada tradicional. Se adopta un procedimiento metodológico trazable: datos abiertos, limpieza reproducible, variables justificadas, etiqueta académica documentada, modelos interpretables y evaluación prudente.

\section{Estrategia general de datos}

La estrategia de datos del TFM se apoya en el Registro de Emergencias del 112 de Castilla y León como fuente empírica principal \cite{Registro112CyL}. Esta decisión aporta coherencia territorial y normativa: el dataset procede de un servicio autonómico 112 real, y el marco de interpretación se apoya en la normativa estatal y de Castilla y León, especialmente la Ley 4/2007, PLANCAL, INFOCAL, INUNCYL y MPCyL. Asimismo, este enfoque garantiza que los escenarios analizados reflejen la casuística real de los incidentes gestionados en la región. De este modo, se evita el uso de datos puramente sintéticos que podrían distorsionar la evaluación del modelo en condiciones operativas.

El conjunto abierto permite trabajar con texto operativo, fecha, localidad, provincia,
coordenadas aproximadas, identificador y categoría preliminar. Sin embargo, no
contiene una prioridad oficial P1--P4 ni todas las variables estructuradas requeridas.
El objetivo es construir una referencia académica mediante supervisión débil y
dejar preparada su revisión humana.

El procesamiento y acondicionamiento de los datos crudos sigue una secuencia estructurada. Este flujo de trabajo garantiza que cada fase se ejecute bajo estrictas salvaguardas metodológicas. En concreto, el flujo de datos se organiza en seis pasos:

\begin{enumerate}
    \item Ingesta y limpieza del Registro 112 CyL.
    \item Normalización de texto, fechas, coordenadas y categorías.
    \item Extracción de señales textuales y variables V01--V15.
    \item Construcción de etiquetas P1--P4 mediante supervisión débil.
    \item Generación de particiones de entrenamiento, validación y prueba, incluida una partición temporal.
    \item Evaluación de Capa 2 y análisis de sesgos, errores y \textit{leakage}.
\end{enumerate}

\section{Registro de Emergencias del 112 de Castilla y León}

El Registro de Emergencias del 112 de Castilla y León es un conjunto de datos público publicado por la Junta de Castilla y León bajo licencia Creative Commons Reconocimiento 4.0 \cite{Registro112CyL}. En este TFM se utiliza el fichero histórico 2008--2022 como fuente real de incidentes para probar el pipeline de datos y alimentar la construcción de etiquetas académicas. Este marco temporal permite abarcar una amplia variedad estacional y tipológica de sucesos históricos registrados.

Su valor principal es el realismo. Los registros no son ejemplos artificiales creados para el TFM, sino incidentes gestionados por un servicio autonómico 112. Esto permite trabajar con lenguaje operativo, categorías heterogéneas, coordenadas, localidades, duplicidades potenciales, ruido textual y limitaciones propias de datos administrativos reales.

\subsection{Campos de entrada}

La Tabla~\ref{tab:campos-registro-cyl-cap7} resume los campos principales del dataset y su relación con los contratos del sistema. A través de esta estructura, se define qué variables del registro bruto se descartan por razones de seguridad o leakage y cuáles se preservan. Este mapeo sistemático constituye el paso previo fundamental para la definición de la interfaz de entrada del motor tri-capa.

\begin{table}[htbp]
\centering
\small
\begin{tabular}{|p{3.2cm}|p{4.6cm}|p{5.4cm}|}
\hline
\textbf{Campo original} & \textbf{Uso en el sistema} & \textbf{Observación metodológica} \\ \hline
\texttt{Titulo} & \texttt{texto\_titulo} en \texttt{IncidentInput} & Fuente textual breve para señales operativas \\ \hline
\texttt{DescripcionBlob} & \texttt{texto\_descripcion} & Requiere limpieza de HTML y normalización \\ \hline
\texttt{FechaIncidente} & \texttt{fecha\_incidente} & Permite análisis temporal y split por año \\ \hline
\texttt{Localidad} & \texttt{localidad} & Contexto territorial no siempre equivalente a coordenada exacta \\ \hline
\texttt{Localidad.1} & \texttt{latitud}, \texttt{longitud} & Permite disponibilidad de coordenadas cuando el campo es válido \\ \hline
\texttt{TipoIncidente} & \texttt{categoria\_preliminar} & Campo orientativo; no basta para priorizar \\ \hline
\texttt{Identificador} & Trazabilidad interna & No se usa como feature predictiva \\ \hline
\texttt{MediosMov} & Solo evaluación ex post si procede & Variable prohibida como entrada por \textit{leakage} \\ \hline
\texttt{PacientesAten} & Solo evaluación ex post si procede & Resultado posterior, no disponible en primera decisión \\ \hline
\texttt{IncidenteCerrado} & Exclusión & Estado final del incidente \\ \hline
\texttt{ultimaActualizacion} & Exclusión & Información posterior al inicio \\ \hline
\end{tabular}
\caption{Campos principales del Registro 112 CyL y uso previsto en el sistema.}
\label{tab:campos-registro-cyl-cap7}
\end{table}

\subsection{Auditoría inicial}

La auditoría del dataset documenta volumen, duplicados, nulos, cobertura temporal, cobertura geográfica y distribución por provincia y categoría. En la versión v0.1.0 esta tarea queda materializada en \texttt{scripts/audit\_dataset.py} y en los informes de \texttt{resources/dataset/audit}. El dataset final utilizado en v0.1.0 conserva 9.380 registros válidos, 41 columnas, cobertura temporal entre 2008-07-28 y 2022-12-31, ausencia de fechas inválidas, 0 identificadores duplicados y cobertura de coordenadas del 99,24\%. Estos datos constituyen la base reproducible del entrenamiento y evaluación, siempre diferenciando los datos reales del Registro 112 CyL de la etiqueta académica P1--P4 generada mediante supervisión débil.

\begin{table}[htbp]
\centering
\small
\begin{tabular}{|p{4.4cm}|p{8.4cm}|}
\hline
\textbf{Elemento auditado} & \textbf{Finalidad} \\ \hline
Número de registros & Confirmar volumen útil para entrenamiento y evaluación \\ \hline
Rango temporal & Verificar coherencia 2008--2022 y detectar fechas anómalas \\ \hline
Distribución por año & Identificar cambios de registro o sesgos temporales \\ \hline
Distribución por provincia & Evitar que la evaluación quede dominada por una provincia \\ \hline
Distribución por \texttt{TipoIncidente} & Conocer la heterogeneidad operativa del dataset \\ \hline
Nulos por columna & Separar campos explotables de campos incompletos \\ \hline
Duplicados o registros equivalentes & Evitar fuga entre train/test por incidentes repetidos \\ \hline
Cobertura de coordenadas & Medir disponibilidad de localización para variables contextuales \\ \hline
Campos ex post & Identificar columnas prohibidas para evitar \textit{leakage} \\ \hline
\end{tabular}
\caption{Elementos obligatorios de la auditoría del dataset.}
\label{tab:auditoria-dataset-cap7}
\end{table}

La auditoría precede a la presentación de resultados para que las métricas puedan interpretarse a partir de la calidad y las limitaciones de los datos.

\section{Limpieza y normalización}

La limpieza transforma el fichero bruto en una representación apta para contratos. Las operaciones implementadas son: lectura robusta del CSV histórico; eliminación de columnas vacías o irrelevantes; normalización de nombres; limpieza de HTML; unión de título y descripción en un texto operativo; normalización para extracción de patrones; parseo de fechas; separación de coordenadas; derivación de provincia; mapeo de \texttt{TipoIncidente} a categoría preliminar; e identificación de columnas prohibidas.

El resultado debe conservar trazabilidad con el identificador original. Cada transformación relevante debe ser reproducible mediante script y, cuando afecte al entrenamiento, quedar versionada. Este criterio de versionado sistemático es el que garantiza la reproducibilidad de todo el pipeline de preparación de datos.

\section{Variables V01--V15}

El sistema no utiliza directamente todas las columnas del dataset. La Capa 1 transforma la entrada en \texttt{IncidentFeatures}, un contrato que contiene variables V01--V15 y señales textuales. La Tabla~\ref{tab:variables-cap7} resume el origen esperado de cada variable.

\begin{table}[htbp]
\centering
\scriptsize
\begin{tabular}{|p{1.0cm}|p{4.0cm}|p{3.9cm}|p{4.0cm}|}
\hline
\textbf{Var.} & \textbf{Variable} & \textbf{Origen principal} & \textbf{Estado v0.1.0} \\ \hline
V01 & Riesgo vital & Texto + señales + etiquetado & Activa \\ \hline
V02 & Número de víctimas estimado & Texto + etiquetado & Activa \\ \hline
V03 & Gravedad de lesiones & Texto + etiquetado & Activa \\ \hline
V04 & Tipo de incidente normalizado & Título, descripción y categoría preliminar & Activa \\ \hline
V05 & Población vulnerable & Texto + reglas + revisión & Activa \\ \hline
V06 & Número de llamadas & Texto, agrupación o indicios de varios avisos & Activa \\ \hline
V07 & Emplazamiento crítico & Texto/localización + reglas & Activa \\ \hline
V08 & Aviso AEMET & Fuente externa & Diferida v0.2.0 \\ \hline
V09 & Condición meteorológica adversa & Texto/fuente externa & Diferida v0.2.0 \\ \hline
V10 & Zona inundable & Cruce geoespacial & Diferida v0.2.0 \\ \hline
V11 & Proximidad Seveso & Registro externo & Diferida v0.2.0 \\ \hline
V12 & Riesgo de propagación & Texto + INFOCAL & Activa \\ \hline
V13 & Multirriesgo & Texto + concurrencia de señales & Activa \\ \hline
V14 & Avisos simultáneos en zona & Texto/agrupación temporal-geográfica & Activa \\ \hline
V15 & Accesibilidad de recursos & Localización + reglas aproximadas & Activa \\ \hline
\end{tabular}
\caption{Variables V01--V15 y estado dentro del alcance v0.1.0.}
\label{tab:variables-cap7}
\end{table}

Las variables V08--V11 se mantienen en el contrato para preservar la evolución futura del sistema, pero no se emplean como entrada del modelo v0.1.0 ni condicionan sus resultados.

\section{Señales operativas textuales}

Las señales textuales son indicadores derivados del título y la descripción. No son etiquetas oficiales, pero ayudan a alimentar tanto la Capa 1 como los anotadores débiles de la Capa 2. Las principales señales son \texttt{signal\_fallecido}, \texttt{signal\_herido\_grave}, \texttt{signal\_atrapado}, \texttt{signal\_intoxicacion}, \texttt{signal\_varias\_llamadas}, \texttt{signal\_incendio}, \texttt{signal\_accidente\_trafico}, \texttt{signal\_rescate}, \texttt{signal\_meteo\_inundacion} y \texttt{riesgo\_vital\_textual}.

Estas señales deben tratarse como indicios. La palabra ``herido'', por ejemplo, no equivale siempre a P1; la expresión ``herido grave atrapado'' sí tiene un peso operativo mucho mayor. Por eso la extracción textual debe combinar patrones, intensificadores y validación posterior, evitando que una búsqueda literal determine por sí sola la prioridad.

\section{Etiqueta académica P1--P4}

La etiqueta P1--P4 es una construcción académica del TFM. No existe como campo oficial en el Registro 112 CyL y no debe presentarse como prioridad institucional. Su función es permitir entrenamiento, comparación y evaluación del sistema.

La guía de etiquetado define cuatro niveles: P1 crítica, P2 alta, P3 media y P4 baja. La etiqueta se asigna en el momento teórico de la primera decisión. Por tanto, no puede utilizar información posterior como medios movilizados, pacientes atendidos o cierre del incidente. La incertidumbre se refleja en la confianza, no rebajando automáticamente la prioridad.

\section{Supervisión débil}

La supervisión débil permite construir etiquetas de entrenamiento a partir de varias funciones de etiquetado imperfectas. En v0.1.0 se implementan cuatro aproximaciones deterministas:

\begin{enumerate}
    \item \textbf{Reglas heurísticas}: aplica una puntuación de gravedad sobre señales textuales, categoría preliminar y variables permitidas.
    \item \textbf{Señales e intensificadores}: aproxima el reconocimiento de entidades nombradas (\textit{Named Entity Recognition}, NER) mediante patrones de gravedad, cantidad y urgencia, sin ejecutar un modelo NER entrenado.
    \item \textbf{Proxy de agrupación semántica}: reutiliza en esta versión una función determinista de puntuación; no ejecuta un algoritmo de agrupamiento vectorial.
    \item \textbf{Proxy de anotador LLM}: incorpora patrones adicionales de vulnerabilidad, emplazamiento y negación, pero no invoca un LLM durante el etiquetado.
\end{enumerate}

Los votos se agregan mediante votación ponderada. Los pesos son \(0{,}85\) para reglas heurísticas, \(1{,}00\) para señales e intensificadores, \(0{,}90\) para el proxy de agrupación y \(1{,}10\) para el proxy de anotador LLM. Se suman los pesos recibidos por cada prioridad y, si existe empate, se escoge la prioridad más urgente. La salida conserva los votos, la etiqueta final, la confianza derivada del peso ganador y el acuerdo local.

\subsection{Acuerdo y control de calidad}

El objetivo mínimo es Krippendorff \(\alpha \geq 0{,}67\). La ejecución canónica alcanza \(\alpha = 0{,}899902\), por encima del umbral establecido. Este valor aporta evidencia de consistencia entre funciones, pero debe interpretarse con cautela: las fuentes comparten patrones y dos de ellas emplean una puntuación muy similar, por lo que no constituyen anotadores plenamente independientes.

Además, se ha preparado una muestra de revisión interna superior al mínimo previsto. Sus campos de decisión humana no están cumplimentados en v0.1.0, por lo que constituye un instrumento preparado y no una validación humana completada.

\subsection{Ablación anti-circularidad}

La supervisión débil introduce el riesgo de que el modelo reproduzca los patrones empleados para construir la propia etiqueta. La ablación sin la fuente de reglas heurísticas mantiene \(\alpha = 0{,}834155\), lo que demuestra que el acuerdo no desaparece al retirar esa fuente. Sin embargo, no garantiza independencia entre la etiqueta y las variables predictivas, porque las funciones restantes comparten señales y criterios con las entradas de Capa 2. Esta dependencia constituye una limitación metodológica de v0.1.0.

\section{Prevención de leakage}

El \textit{leakage} es uno de los riesgos metodológicos más graves del proyecto. Se produce cuando el modelo utiliza información que no estaría disponible en el momento de la primera decisión. En el Registro 112 CyL existen campos que pueden reflejar actuaciones posteriores o resultados finales. La Tabla~\ref{tab:leakage-cap7} resume las columnas prohibidas.

El sistema excluye de forma explícita variables que no estarían disponibles en el momento de la primera decisión, como recursos movilizados, pacientes atendidos, estado de cierre o información posterior al desarrollo del incidente. Esta decisión evita \textit{leakage} y mantiene la evaluación alineada con el objetivo de apoyo a la primera decisión. Al asegurar que el modelo solo utiliza datos pre-decisionales, se simula fielmente el entorno real de un despachador del 112. Esto garantiza que las métricas de precisión obtenidas sean realistas y extrapolables a la práctica operativa.

\begin{table}[htbp]
\centering
\small
\begin{tabular}{|p{4.2cm}|p{8.2cm}|}
\hline
\textbf{Columna} & \textbf{Razón de exclusión como feature} \\ \hline
\texttt{MediosMov} & Refleja recursos movilizados, es decir, una consecuencia de la decisión operativa \\ \hline
\texttt{medios\_mov\_limpio} & Derivado de \texttt{MediosMov}; mantiene la misma fuga \\ \hline
\texttt{medios\_mov\_uso\_ recomendado} & Codificación auxiliar ex post; no disponible en primera decisión \\ \hline
\texttt{PacientesAten} & Resultado clínico o asistencia posterior \\ \hline
\texttt{pacientes\_aten\_limpio} & Derivado de \texttt{PacientesAten}; conserva información posterior \\ \hline
\texttt{IncidenteCerrado} & Estado final del incidente \\ \hline
\texttt{ultimaActualizacion} & Actualización posterior al registro inicial \\ \hline
\texttt{Enlace al contenido} & Identificador externo; útil para trazabilidad, no como predictor \\ \hline
\texttt{Unnamed: 13} & Columna vacía o no informativa \\ \hline
\end{tabular}
\caption{Columnas prohibidas como entrada del motor por riesgo de \textit{leakage}.}
\label{tab:leakage-cap7}
\end{table}

Estas columnas pueden conservarse en zonas controladas de auditoría o evaluación ex post, siempre separadas del entrenamiento e inferencia. El repositorio incluye un test de leakage que escanea los módulos de la Capa 1 y la Capa 2 para impedir referencias no permitidas. Dicho test actúa como una salvaguarda de integración continua en el ciclo de desarrollo de software. Gracias a esta verificación automática, se garantiza que ningún desarrollador introduzca inadvertidamente variables ex post en los estimadores.

\begin{table}[htbp]
\centering
\scriptsize
\begin{tabular}{|p{3.4cm}|p{4.6cm}|p{5.1cm}|}
\hline
\textbf{Grupo} & \textbf{Columnas o variables} & \textbf{Uso permitido} \\ \hline
Texto inicial & \texttt{Titulo}, \texttt{DescripcionBlob}, texto operativo limpio & Entrada para la Capa 1 y extracción de señales \\ \hline
Contexto predecisional & Fecha, localidad, provincia, coordenadas, \texttt{TipoIncidente} & Normalización, taxonomía, análisis temporal y territorial \\ \hline
Variables derivadas permitidas & V01--V07, V12--V15 y señales textuales & Entrada de la Capa 2 y explicabilidad \\ \hline
Identificadores & \texttt{Identificador} & Trazabilidad, splits y auditoría; no como predictor \\ \hline
Columnas prohibidas & \texttt{MediosMov}, \texttt{PacientesAten}, \texttt{IncidenteCerrado}, \texttt{ultimaActualizacion} y derivados & Solo auditoría o análisis ex post; nunca entrenamiento ni inferencia \\ \hline
\end{tabular}
\caption{Columnas permitidas y prohibidas según disponibilidad en la primera decisión.}
\label{tab:permitidas-prohibidas-leakage-cap7}
\end{table}

\section{Particiones de entrenamiento, validación y prueba}

Para garantizar la validez de los modelos entrenados y medir su capacidad de generalización, se definen particiones rigurosas que impiden la fuga de información entre conjuntos. La división de datos debe impedir que el modelo vea en entrenamiento información demasiado parecida a la de test, controlando sesgos geográficos y temporales. Con este propósito, se plantean dos estrategias de partición complementarias: un split estratificado para la evaluación general y un split temporal para evaluar la robustez frente al cambio de periodo operativo.

\subsection{Partición estratificada}

La partición principal divide el conjunto en 6.512 registros de entrenamiento, 1.415 de validación y 1.453 de prueba, preservando la distribución por prioridad, año y provincia con semilla fija 42. La partición de prueba contiene P1=655, P2=349, P3=375 y P4=74, conforme al reporte canónico \texttt{artifacts/reports/evaluation\_v0.1.0.json}.

\begin{table}[htbp]
\centering
\small
\begin{tabular}{|p{3.2cm}|p{3.0cm}|p{5.4cm}|}
\hline
\textbf{Partición} & \textbf{Filas} & \textbf{Distribución P1--P4} \\ \hline
Entrenamiento estratificado & 6.512 & P1=3.049; P2=1.504; P3=1.701; P4=258 \\ \hline
Validación estratificada & 1.415 & P1=658; P2=329; P3=374; P4=54 \\ \hline
Test estratificado & 1.453 & P1=655; P2=349; P3=375; P4=74 \\ \hline
Test temporal 2022 & 735 & Evaluación de estabilidad fuera del periodo principal \\ \hline
\end{tabular}
\caption{Distribución de las particiones adoptadas en v0.1.0.}
\label{tab:splits-final-cap7}
\end{table}

\subsection{Partición temporal}

La partición temporal evalúa robustez fuera del periodo de entrenamiento. La versión ejecutada entrena con registros hasta 2020, valida con 2021 y reserva 2022 como prueba temporal. Los tamaños son 7.935, 710 y 735 registros, respectivamente.

Ambos splits son necesarios. El estratificado mide rendimiento general controlando distribución; el temporal mide estabilidad ante cambio de periodo. La combinación de ambas perspectivas metodológicas proporciona una imagen completa del comportamiento del sistema. Mientras que el primero valida el acierto en condiciones estándar, el segundo pone a prueba la resiliencia del modelo frente al paso del tiempo.

\section{Uso en entrenamiento y evaluación}

El dataset resultante alimenta tres usos distintos: pruebas de extracción de la Capa 1, entrenamiento y evaluación de la Capa 2, y evaluación transversal de sesgo, falsos negativos P1, calibración y matriz de confusión. La Capa 3 no se entrena con el dataset operativo; utiliza las salidas de la Capa 2 y el corpus normativo para generar explicaciones. Cada uno de estos usos cuenta con sus propios protocolos de verificación y contratos de entrada. De esta forma, se mantiene una separación clara de responsabilidades entre los distintos componentes del sistema.

\section{Limitaciones del dataset}

Las limitaciones deben quedar explicitadas para evitar sobreprometer. La primera es la ausencia de prioridad oficial. Los resultados obtenidos deben interpretarse en el marco de etiquetas débiles: la etiqueta P1--P4 no procede del servicio 112 ni de una prioridad oficial, sino de una construcción académica trazable y reproducible. Cualquier cambio en la guía de etiquetado, en los splits o en el dataset limpio obliga a repetir la evaluación.

La segunda limitación es la granularidad limitada: el dataset contiene registros administrativos, no la secuencia completa de llamadas minuto a minuto. La tercera es que muchas variables deben inferirse del texto. La cuarta es que AEMET, inundabilidad y Seveso quedan fuera del núcleo v0.1.0. La quinta es que la validación externa con personal operativo del 112 queda como trabajo futuro. La sexta es el posible sesgo territorial, temporal o por categoría.

\section{Conclusiones}

El diseño de datos presentado en este capítulo conecta el Registro 112 CyL con la arquitectura de tres capas del sistema. La estrategia evita dos extremos: no inventa un dataset totalmente sintético, pero tampoco finge disponer de una prioridad oficial que no existe. Esta postura pragmática asegura que el desarrollo técnico esté cimentado sobre bases realistas y éticas. Como consecuencia, las conclusiones derivadas del comportamiento del modelo resultan defendibles ante evaluadores y profesionales de emergencias.

La aportación metodológica consiste en convertir datos abiertos reales en un conjunto explotable mediante limpieza, contratos, variables operativas, supervisión débil, control de leakage y splits reproducibles. Esta base es la que permite que la Capa 2 pueda entrenarse y evaluarse de forma defendible. Además, este marco conceptual establece un precedente para la reutilización de otros registros administrativos de la administración pública. Al estructurar metodológicamente datos crudos y ruidosos, se abre la puerta a su aprovechamiento en tareas complejas de toma de decisiones.

El siguiente capítulo documenta la implementación del prototipo software que ejecuta este flujo: backend, orquestador, herramientas MCP, logging, modo degradado e interfaz mínima. Se detallará la arquitectura del código, las dependencias tecnológicas seleccionadas y las decisiones de diseño del backend. De igual manera, se describirán las pruebas de integración que aseguran el correcto intercambio de información entre las tres capas.
